\section{Sketches and future work}\label{sec:sketches}

This chapter records candidates, not laws.  Several structural
candidates were considered for the catalog but are not laws of it,
because their recognition source has not been identified.  A
recognition source is the substrate, apparatus layer or carrier
discipline that makes a named structure precise enough to state in the
template of \Cref{def:normal-form}: a setup, a theorem and a proof.
Without one, a phrase may be suggestive and still fail to determine
these.  Some of the candidates below may turn out to belong to a
substrate paper only, because the structure is carrier-specific; others
may become laws once a recognition source is supplied.

\paragraph{Dimension as stable rank of access.}  The sketch asks
whether dimension can be read, at the layer-agnostic level, as a
presentation-invariant rank of the access quotient: roughly, the stable
number of independent directions that survive the declared observation
discipline.  Such a reading would explain why dimension is neither raw
carrier data nor a coordinate convention.  To state it as a law one
must know which rank notion is meant, which class of quotient
presentations preserves it, and which audit certifies that the rank is
observed rather than imposed.  These data are not part of the present
apparatus.

\paragraph{Smoothness as coherent linearization.}  The sketch asks
whether smooth structure can be characterized by compatible local
linearizations across declared overlaps: local linear approximations
should globalize exactly when they descend through the relevant access
and overlap quotients.  This would connect the local--global
obstruction of \Cref{sec:transport:local-global} with a
smoothness-specific substrate, but it needs a substrate-independent
account of what counts as a linearization, which overlaps preserve it,
and how the associated coherence ledger is audited.

\paragraph{Stability of nothing.}  A ``nothing'' state at a layer is
rarely absolute absence; it is more naturally a terminal, initial, null
or background class relative to a quotient and an audit.  The catalog
already contains Vacuum / Background Selection (\Cref{thm:F50}): a
vacuum is the single element of the declared vacuum criterion, with
formed target.  The sketch asks the stronger question of when the
stability of the null class forces a nonempty residual structure, which
needs a recognition source for the null class and for the residual
supposed to survive it.

\paragraph{Other candidates.}  Boundary compression, topology,
continuum behavior, critical loci, uncertainty and inaccessible
determining state are treated by laws of the catalog: Horizon Sufficiency, Information Loss, Continuum Emergence,
Criticality, Complementarity as Non-Joint Access, Topology as Global
Gluing Invariant and Hiddenness.  Other formulations on these six
topics are not stated here; in the application of \Cref{thm:F52} below
they count as named under those topics.

\paragraph{Scope of the catalog.}  The eight axes are an organizing
scheme that makes the catalog navigable, not a structured object, and
the catalog is not claimed to be complete.  The Closure Completeness
Boundary (\Cref{thm:F52}) characterizes scoped completeness for any
declared residuals and directions.  Take as \(D\) the candidate
directions considered for the catalog, let
\(\operatorname{inScope}(d)\) mean that \(d\) is a law of the catalog
or reduces to one, let \(\operatorname{namedOutside}(d)\) mean that
\(d\) is named in this chapter, and let \(\operatorname{claimed}(d)\)
hold for every \(d\in D\).  The direction clause of \Cref{thm:F52} then
requires every candidate to be covered or named.  Since the list of
candidates is not reproduced here, the paper does not verify that
clause; it records only the conditional statement that if every
candidate is a law, reduces to one, or is named here, then the direction
clause holds and the unmarked frontier is empty.  The residual clauses
of \Cref{thm:F52} would require a residual register for the catalog,
which is not given.  The candidates named here that neither belong to
nor reduce to a law form the catalog boundary \(\partial D\) of
\Cref{def:catalog-boundary}.  \Cref{thm:F52} is not an absolute
completeness theorem: a new recognition source may move a sketch out of
\(\partial D\) and into the catalog.

A candidate that comes with a recognition source can be tested in the
same way as the laws of this paper: identify the setup, state the
quotient or closure normal form, give the proof, record the cases and
write the nonclaims.  If it passes without special carrier features, it
belongs in the catalog; if it depends on special analytic, geometric,
arithmetic or physical structure, it belongs in the corresponding
substrate paper.
